\chapter{Los números complejos}
\label{cap:complejos}

\section{Construcción y forma binómica}
\label{sec:complejos-binomica}
El cuerpo $\C$ permite resolver $z^2=d$ también para reales negativos.
Se construye sobre $\R^2$ con las operaciones
\begin{align*}
 (a,b)+(c,d)&=(a+c,b+d),\\
 (a,b)(c,d)&=(ac-bd,ad+bc).
\end{align*}
El cero es $(0,0)$, la unidad $(1,0)$ y el opuesto de $(a,b)$ es
$(-a,-b)$. Si $(a,b)\ne(0,0)$, su inverso es
\[
 (a,b)^{-1}=\left(\frac a{a^2+b^2},\frac{-b}{a^2+b^2}\right).
\]
Estas operaciones hacen de $\C$ un cuerpo. La aplicación
$a\mapsto(a,0)$ incluye $\R$ como subcuerpo. Definiendo $i=(0,1)$,
se tiene $i^2=-1$, y cada complejo se escribe de forma única como
\[
 z=a+bi\qquad(a,b\in\R).
\]
Esta es la \textbf{forma binómica}. La parte real es $\operatorname{Re}z=a$
y la parte imaginaria es $\operatorname{Im}z=b$, que es un número real,
no el término $bi$. El complejo es real si $b=0$ y es imaginario puro
si $a=0$. La igualdad de complejos exige igualdad de ambas partes.

Para sumar se suman las partes correspondientes; para multiplicar se
distribuye y se sustituye $i^2$ por $-1$.
Las potencias de $i$ se repiten cada cuatro exponentes:
\[
 i^0=1,\quad i^1=i,\quad i^2=-1,\quad i^3=-i,\quad i^{n+4}=i^n.
\]
Para un exponente entero basta conocer su resto módulo $4$.

\section{Conjugación, módulo y plano complejo}
\label{sec:conjugacion-modulo}
El \textbf{conjugado} de $z=a+bi$ es $\bar z=a-bi$. Se cumple
\[
 \overline{\bar z}=z,\qquad
 \overline{z+w}=\bar z+\bar w,\qquad
 \overline{zw}=\bar z\bar w,
 \qquad z\bar z=a^2+b^2.
\]
También $z+\bar z=2\operatorname{Re}z$ y
$z-\bar z=2i\operatorname{Im}z$. Por tanto,
$z\in\R\eqv\bar z=z$ y $z$ es imaginario puro si y solo si
$\bar z=-z$. La conjugación conserva potencias y, para denominadores
no nulos, inversos y cocientes.

El \textbf{módulo} es $|z|=\sqrt{z\bar z}=\sqrt{a^2+b^2}$.
Al representar $z$ por el punto $(a,b)$ del plano, llamado su punto
de afijo $z$, el módulo es la distancia al origen. Se tiene
\begin{align*}
 |z|&\ge0,& |z|=0&\eqv z=0,\\
 |zw|&=|z||w|,&|\bar z|&=|z|,\\
 |z+w|&\le|z|+|w|,&\bigl||z|-|w|\bigr|&\le|z-w|.
\end{align*}
La suma corresponde a la suma de vectores; el conjugado refleja el punto
respecto al eje real y el opuesto respecto al origen.

El conjugado facilita la división: si $w\ne0$,
\[
 \frac1w=\frac{\bar w}{|w|^2},\qquad
 \frac zw=\frac{z\bar w}{|w|^2}.
\]
Así se transforma el denominador en un real positivo antes de separar
las partes real e imaginaria.

\section{Argumento y formas trigonométrica, polar y exponencial}
\label{sec:complejos-polar}
Para $z\ne0$, su \textbf{argumento} es el ángulo orientado desde el eje
real positivo al vector de afijo $z$. Se determina módulo $2\pi$:
si $\alpha$ es un argumento, todos son $\alpha+2k\pi$, $k\in\Z$.
El cero no tiene argumento.

Si $r=|z|>0$ y $\alpha$ es un argumento,
\[
 z=r(\cos\alpha+i\sin\alpha),\qquad
 a=r\cos\alpha,\quad b=r\sin\alpha.
\]
La primera expresión es la \textbf{forma trigonométrica}; la notación
$r_\alpha$ es la \textbf{forma polar}. Para recuperar el argumento desde
$a,b$ hay que respetar el cuadrante: el cociente $b/a$, cuando está
definido, no determina por sí solo el ángulo módulo $2\pi$.

Se usa la notación $e^{i\alpha}=\cos\alpha+i\sin\alpha$, con la cual
$z=re^{i\alpha}$ es la \textbf{forma exponencial}. Para $r,s>0$,
\begin{align*}
 re^{i\alpha}=se^{i\beta}&\eqv r=s\ \text{y}\ \alpha\equiv\beta\pmod{2\pi},\\
 \overline{re^{i\alpha}}&=re^{-i\alpha},\qquad
 -re^{i\alpha}=re^{i(\alpha+\pi)},\\
 (re^{i\alpha})(se^{i\beta})&=rs\,e^{i(\alpha+\beta)},\\
 \frac{re^{i\alpha}}{se^{i\beta}}&=\frac rs e^{i(\alpha-\beta)},\\
 (re^{i\alpha})^n&=r^ne^{in\alpha}\qquad(n\in\Z).
\end{align*}
El producto multiplica módulos y suma argumentos; el cociente divide
módulos y resta argumentos. Multiplicar por un complejo no nulo realiza
una rotación y una homotecia de razón positiva respecto al origen.
Estas formas son útiles para productos, cocientes y potencias; para
sumar suele convenir volver a la forma binómica.

La \textbf{fórmula de De Moivre} es
\[
 (\cos\alpha+i\sin\alpha)^n=\cos(n\alpha)+i\sin(n\alpha)
 \qquad(n\in\Z).
\]
Al compararla con el binomio de Newton se relacionan las funciones
trigonométricas de múltiplos de un ángulo. Las \textbf{fórmulas de Euler}
son
\[
 \cos\alpha=\frac{e^{i\alpha}+e^{-i\alpha}}2,\qquad
 \sin\alpha=\frac{e^{i\alpha}-e^{-i\alpha}}{2i}.
\]

\section{Raíces y ecuaciones de segundo grado}
\label{sec:raices-complejas}
Sea $w=\rho e^{i\beta}\ne0$ y $n\in\N^*$. Las soluciones de $z^n=w$ son
\[
 z_k=\sqrt[n]{\rho}\,
 e^{i(\beta+2k\pi)/n},\qquad k=0,1,\ldots,n-1.
\]
Hay exactamente $n$ raíces distintas: todas tienen módulo
$\sqrt[n]\rho$ y sus argumentos están separados por $2\pi/n$.
Para $n\ge3$ son los vértices de un polígono regular centrado en el
origen. Cambiar el argumento elegido para $w$ solo permuta las raíces.
Si $w=0$, la única raíz es $0$.

Las \textbf{raíces $n$-ésimas de la unidad} son
$e^{2k\pi i/n}$, $0\le k<n$.
La expresión ``una raíz'' de un complejo debe distinguirse del conjunto
de todas sus raíces. La regla real $\sqrt{xy}=\sqrt x\sqrt y$, para
$x,y\ge0$, no se extiende sin más a raíces elegidas de complejos.

Para hallar raíces cuadradas en forma binómica de $w=u+vi$, se busca
$z=x+yi$ y se usa
\[
 x^2-y^2=u,\quad 2xy=v,\quad x^2+y^2=|w|.
\]
De aquí $x^2=(|w|+u)/2$ e $y^2=(|w|-u)/2$; los signos deben satisfacer
$2xy=v$. Si $w\ne0$, las dos raíces son opuestas.

Para $az^2+bz+c=0$ con $a,b,c\in\C$ y $a\ne0$:
\begin{enumerate}
\item Calcular $\Delta=b^2-4ac$.
\item Obtener una raíz cuadrada $\varepsilon$ de $\Delta$, por forma polar
o por el sistema binómico anterior. Si $\Delta\ge0$ es real, puede
tomarse $\sqrt\Delta$; si es real negativo, $i\sqrt{-\Delta}$.
\item Las soluciones son
$\displaystyle z=\frac{-b+\varepsilon}{2a}$ y
$\displaystyle z=\frac{-b-\varepsilon}{2a}$.
\end{enumerate}
Son distintas si $\Delta\ne0$; si $\Delta=0$, coinciden en una raíz doble.

\section{Aplicaciones geométricas}
\label{sec:geometria-compleja}
Sean $z_A,z_B,z_C$ los afijos de puntos del plano.
El vector $\overrightarrow{AB}$ tiene afijo $z_B-z_A$; la distancia
entre $A,B$ es $|z_B-z_A|$ y el punto medio tiene afijo
$(z_A+z_B)/2$.

Si $B\ne A$ y $C\ne A$, el ángulo orientado de
$\overrightarrow{AB}$ a $\overrightarrow{AC}$ es
\[
 \arg\left(\frac{z_C-z_A}{z_B-z_A}\right)\pmod{2\pi}.
\]
Los tres puntos están alineados si y solo si ese cociente es real;
las rectas $AB,AC$ son perpendiculares si y solo si es imaginario puro.
Las condiciones de no coincidencia aseguran que ambos vectores sean
no nulos y que el cociente tenga argumento.

Las transformaciones se expresan directamente sobre afijos:
\begin{center}
\begin{tabular}{@{}ll@{}}
\toprule
Transformación & Afijo transformado $z'$\\
\midrule
Traslación de vector de afijo $w$ & $z+w$\\
Rotación de centro $w$ y ángulo $\alpha$ & $w+e^{i\alpha}(z-w)$\\
Homotecia de centro $w$, razón $k\in\R^*$ & $w+k(z-w)$\\
Simetría respecto al eje real & $\bar z$\\
Simetría en recta por $0$ de dirección $\alpha$ & $e^{2i\alpha}\bar z$\\
\bottomrule
\end{tabular}
\end{center}
Para actuar respecto a un centro $w$, se traslada primero al origen,
se aplica la transformación y se deshace la traslación.
La simetría en una recta por el origen se obtiene girando esa recta al
eje real, conjugando y deshaciendo el giro.

\section{Orden y teorema fundamental del álgebra}
No existe un orden total en $\C$ compatible con sus operaciones de
cuerpo. En un cuerpo ordenado todos los cuadrados son no negativos,
pero $i^2=-1$ lo impediría, porque $1>0$.
El módulo permite comparar distancias, sin proporcionar ese orden algebraico.

El \textbf{teorema fundamental del álgebra} afirma que todo polinomio
de coeficientes complejos y grado al menos uno tiene una raíz compleja.
Como consecuencia, si su grado es $n\ge1$,
\[
 p(z)=a(z-z_1)\cdots(z-z_n),\qquad a\ne0.
\]
Las raíces $z_1,\ldots,z_n$ pueden repetirse: hay $n$ raíces contando
multiplicidades, aunque no necesariamente $n$ valores distintos.
